\section{RISK-AWARE RECOVERY CONSTRAINT SYNTHESIS}
\label{sec-risk-synthesis}
The belief graph in Section~\ref{sec-patbg} limits how adversarial evidence changes the online knowledge state, but a bounded belief error can still be harmful if it is converted directly into field authority. This section therefore constructs an explicit security boundary between uncertain recovery knowledge and executable control. The design realizes the admissibility model in Section~\ref{subsec-risk-admissibility}: evidence-dependent prerequisites are checked probabilistically, invariant engineering constraints remain authoritative, and learned feasibility models are used only to reduce verification cost.

\subsection{Belief-Gated Semantic Prerequisites}
For candidate action $\mathbf{a}$, let $\mathcal{Q}(\mathbf{a})$ denote the runtime assertions required for execution. Typical examples include controller availability, command-path reachability, measurement trust, authorization state, and protection readiness. We define the semantic risk estimate
\begin{align}
\label{eq-action-semantic-risk}
r_B(\mathbf{a},t)=\sum_{q\in\mathcal{Q}(\mathbf{a})}\delta_q,
\end{align}
where each prerequisite must satisfy
\begin{align}
\label{eq-action-belief-test}
b_t(q)\ge 1-\delta_q
\end{align}
and $r_B(\mathbf{a},t)\le\bar\delta_B$. The parameters $\delta_q$ are operation-dependent risk budgets rather than arbitrary confidence thresholds. Safety-critical switching or protection actions can use smaller budgets than reversible communication rerouting or low-impact advisory actions.

An assertion whose belief is below its required threshold is not automatically declared false; instead the associated action is withheld until the prerequisite is independently verified or a different action not relying on that assertion is selected. This distinction prevents an attacker from turning evidence suppression into a fabricated physical fact while preserving the principle that uncertainty must not create control authority.

For each semantically admitted action, the controller generates an authorization context
\begin{align}
\label{eq-authorization-context}
\mathcal{C}_t(\mathbf{a})=
\big\{\mathcal{Q}(\mathbf{a}),\,\{b_t(q)\},\,\mathcal{P}_t(\mathbf{a}),\,v_B\big\},
\end{align}
where $\mathcal{P}_t(\mathbf{a})$ contains the provenance groups supporting the prerequisites and $v_B$ is the belief-graph version. This context is retained for post-action audit and for updating source reliability after independently verifiable outcomes become available.

\subsection{Invariant Verification and Candidate Selection}
Let $\mathcal{C}_t^{0}$ be a bounded candidate set generated by decentralized restoration policies or by local action expansion. We first remove candidates that fail the belief gate,
\begin{align}
\label{eq-stage-belief-gate}
\mathcal{C}_t^{B}=\mathcal{C}_t^{0}\cap\mathbb{U}_t^{B}.
\end{align}
A learned surrogate $\widehat v_\chi(\mathbf{a})$ may then estimate physical infeasibility or restoration cost, but its output is not interpreted as a safety certificate. We therefore use it primarily for ranking,
\begin{align}
\label{eq-surrogate-ranking}
\operatorname{score}_t(\mathbf{a})=
\widehat Q_t(\mathbf{a})
-\lambda_v\widehat v_\chi(\mathbf{a})
-\lambda_r r_B(\mathbf{a},t),
\end{align}
where $\widehat Q_t$ is a policy or value estimate. Candidates are checked in descending score order by the authoritative verifier.

The final verifier evaluates invariant constraints over the conservative state set $\widehat{\mathcal X}_t$ from Eq.~\ref{eq-invariant-admissible-set}. Define
\begin{align}
\label{eq-authoritative-verifier}
\mathsf{Verify}(\mathbf{a},\widehat{\mathcal X}_t)=1
\end{align}
only if all required power-balance, voltage, thermal, radiality, device, and protection constraints hold for the verifier's uncertainty model. Depending on the action type, the implementation may call AC power flow, a certified conservative approximation, reachability analysis, or a mixed-integer restoration check. The verifier is authoritative; neither the policy value nor the surrogate can override a rejection.

To avoid committing to a fixed top-$K$ that may reject all safe actions because of surrogate ranking error, verification proceeds in bounded batches. If the highest-ranked batch contains no feasible action, the next batch is checked until a feasible action is found or the real-time verification budget is exhausted. If no nontrivial action passes, a certified fallback $\mathbf{a}^{fb}_t$ is executed.

\begin{algorithm}[t]
\caption{Risk-Aware Recovery Constraint Synthesis}
\label{alg-risk-gate}
\KwIn{Candidate set $\mathcal{C}_t^{0}$; belief graph $\mathcal{G}_t^B$; conservative state set $\widehat{\mathcal X}_t$; risk budgets $\{\delta_q\},\bar\delta_B$; verification budget $K_{max}$}
\KwOut{Executable action $\mathbf{a}^{exec}_t$ and authorization record $\mathcal{R}^{auth}_t$}
$\mathcal{C}_t^{B}\leftarrow\emptyset$\;
\ForEach{$\mathbf{a}\in\mathcal{C}_t^{0}$}{
  Query prerequisite set $\mathcal{Q}(\mathbf{a})$ and current beliefs $\{b_t(q)\}$\;
  \If{$b_t(q)\ge1-\delta_q$ for all $q\in\mathcal{Q}(\mathbf{a})$ and $\sum_q\delta_q\le\bar\delta_B$}{
    Build authorization context $\mathcal{C}_t(\mathbf{a})$\;
    Insert $\mathbf{a}$ into $\mathcal{C}_t^{B}$\;
  }
}
Compute surrogate/value score in Eq.~\ref{eq-surrogate-ranking} for each $\mathbf{a}\in\mathcal{C}_t^{B}$ and sort descending\;
$k\leftarrow0$\;
\ForEach{ranked $\mathbf{a}\in\mathcal{C}_t^{B}$}{
  \If{$k\ge K_{max}$}{break}
  $k\leftarrow k+1$\;
  \If{$\mathsf{Verify}(\mathbf{a},\widehat{\mathcal X}_t)=1$}{
    $\mathbf{a}^{exec}_t\leftarrow\mathbf{a}$\;
    Create $\mathcal{R}^{auth}_t$ containing $v_B$, prerequisite beliefs, provenance groups, and verifier result\;
    \Return{$(\mathbf{a}^{exec}_t,\mathcal{R}^{auth}_t)$}\;
  }
}
$\mathbf{a}^{exec}_t\leftarrow\mathbf{a}^{fb}_t$ and record fallback reason\;
\Return{$(\mathbf{a}^{exec}_t,\mathcal{R}^{auth}_t)$}\;
\end{algorithm}

\subsection{Complexity Analysis}
Let $M_t=|\mathcal{C}_t^{0}|$ be the number of candidate actions, $\bar q$ the maximum number of evidence-dependent prerequisites per action, $C_{sur}$ the cost of one surrogate/value evaluation, and $C_{phy}$ the cost of one authoritative physical verification. With indexed prerequisite relations, semantic queries and belief checks cost $O(M_t\bar q)$. Ranking costs $O(M_t\log M_t)$ in the worst case, or $O(M_t+K_{max}\log M_t)$ with a bounded priority queue. At most $K_{max}$ candidates invoke the expensive physical verifier. Hence Algorithm~\ref{alg-risk-gate} has
\begin{align}
\label{eq-complexity-risk-gate}
T_{\rm gate}=O\!\left(M_t\bar q+M_tC_{sur}+M_t\log M_t+K_{max}C_{phy}\right).
\end{align}
Because $C_{phy}$ typically dominates, the practical objective is to keep $K_{max}\ll M_t$ without granting the surrogate authorization authority. The working memory is $O(M_t\bar q)$ if all authorization contexts are materialized simultaneously and can be reduced to $O(M_t+\bar q)$ by constructing contexts lazily in ranked order. The complexity statement intentionally leaves $C_{phy}$ abstract because its value depends on whether the verifier uses sparse AC power flow, conservative linear screening, reachability, or mixed-integer optimization.
